% ======================================================================
\section{The cosmic mass}\label{sec:cosmic-mass}
% ======================================================================

\emph{What number can carry a mass and a charge at once?} This section proposes one number for each body, fixes the size of its second unit from two measured constants, says which angle the frame is drawn at and why that choice is free, and says how the number is read. The statements that the paper assumes are labelled postulates: the two of this section, and the two laws of \cref{sec:rest,sec:motion}.

\subsection{One number for mass and charge}\label{subsec:cm-def}

\begin{postulate}[Cosmic mass]\label{def:cosmic-mass}
The \emph{cosmic mass} of a body of mass $m$ and electric charge $q$ is the General Number
\[
\mu = m + q\,\Gp ,
\]
in kilograms. Its real part is the mass $m$, its Composer coordinate is the charge $q$, and its Composer part $q\,\Gp$, the \emph{charge part}, is in kilograms, so that the Composer Unit $\Gp$ carries the unit kg\,C$^{-1}$.
\end{postulate}

\begin{figure}[t]
\centering
% The cosmic mass as a General Number (Postulate 4.1). Drawn at 55 degrees, not 79.86, so that the lean is visible.
\begin{tikzpicture}[>=Stealth, scale=1.05, every node/.style={font=\small}]
  \def\ang{55}
  \begin{scope}
    \clip (-0.5,-0.6) rectangle (7.3,3.5);
    \foreach \a in {-3,...,7} \draw[black!12] ($(\a,0)+(\ang:-2)$) -- ($(\a,0)+(\ang:6)$);
    \foreach \b in {-1,...,5} \draw[black!12] ($(-3,0)+(\ang:{\b*0.8})$) -- ($(9,0)+(\ang:{\b*0.8})$);
  \end{scope}
  \draw[->, black!55] (-0.5,0) -- (7.3,0) node[below left] {mass axis (real)};
  \draw[->, black!55] (0,0) -- (\ang:4.1) node[above] {charge axis (Composer)};
  \draw[->, very thick, black!80] (0,0) -- (3.2,0) node[midway, below] {$m$};
  \draw[->, very thick, blue!60!black] (3.2,0) -- ($(3.2,0)+(\ang:2.4)$) node[midway, right] {$q\,\Gp$};
  \draw[->, thick, densely dashed] (0,0) -- ($(3.2,0)+(\ang:2.4)$);
  \fill ($(3.2,0)+(\ang:2.4)$) circle (1.5pt) node[above right] {$\mu = m + q\,\Gp$};
  \draw[black!55] (0.7,0) arc (0:\ang:0.7);
  \node[black!65] at ({\ang/2}:1.0) {$\pb$};
  \fill (0,0) circle (1.2pt) node[below left] {$0$};
\end{tikzpicture}

\caption{The cosmic mass of a body as a General Number (\cref{def:cosmic-mass}): $m$ steps along the mass axis and $q$ steps of the Composer Unit along the charge axis. The charge axis is drawn at $55^\circ$ and the two parts with comparable lengths, so that the construction can be seen; the paper's frame is drawn at $79.86^\circ$, and for a charged particle the charge part is about $10^{18}$ to $10^{21}$ times longer than the mass.}
\label{fig:cosmic-mass}
\end{figure}

\Cref{fig:cosmic-mass} draws the construction. A neutral body is a real cosmic mass. The free neutron, for example, is the real number $\mu_n = m_n$. A body with charge has a second part, and for the elementary particles that part dominates. With the size of the Composer Unit found below, the charge part of the electron, $e\,\lvert\Gp\rvert \approx 1.85921 \times 10^{-9}$~kg, is about $2.04 \times 10^{21}$ times its mass, and that of the proton about $1.11 \times 10^{18}$ times its mass.

\subsection{The size and the angle of the Composer Unit}\label{subsec:our-plane}

The Composer Unit has a size and an angle. The size is a physical quantity; the angle, as we shall see, belongs to the frame in which the plane is drawn. The size is fixed by asking that the term of two charges in the product of two cosmic masses (\cref{sec:rest}) have the strength of Coulomb's law:
\begin{equation}\label{eq:cu-size}
G\,s^2 = K , \qquad s = \lvert\Gp\rvert = \sqrt{\frac{K}{G}} = c\,\sqrt{\frac{\mu_0}{4\pi G}} \approx 1.16043 \times 10^{10}\ \mathrm{kg\,C^{-1}} ,
\end{equation}
where the second form uses $K = \mu_0 c^2/(4\pi)$, with $\mu_0$ the magnetic constant. This is Stoney's ratio of mass to charge~\citep{Stoney1881}. With the CODATA 2022 values~\citep{CODATA2022} it is known to about one part in $10^5$, half the relative uncertainty of $G$. The constant $K$ therefore enters the framework as an input; the Coulomb term has Coulomb's strength because $s$ was chosen to give it.

The Axiom~\cref{eq:axiom} compares the length of the Composer Unit with the pure number $\lvert\cos\pb + \sin\pb\rvert$, while $s$ carries the unit kg\,C$^{-1}$. To give it an angle, a scale unit $u$ must first be chosen, and the Axiom is read for the pure number $A = s/u$. With $u = 10^{10}$~kg\,C$^{-1}$, $A = 1.16043$ and
\begin{equation}\label{eq:our-angle}
\lvert\cos\pb + \sin\pb\rvert = A = 1.16043 .
\end{equation}
We keep the two symbols apart: $s$ is the physical size, fixed by the measured constants $G$ and $K$, and $A$ is the drawing length, fixed by the choice of $u$.
By \citep[Theorem~3.6]{DerrazGN2026}, two angles between $0$ and $\pi/2$ satisfy~\cref{eq:our-angle}:
\[
\pb_1 = \arcsin\Bigl(\frac{1.16043}{\sqrt{2}}\Bigr) - \frac{\pi}{4} \approx 10.14^\circ , \qquad \pb_2 = \frac{\pi}{2} - \pb_1 \approx 79.86^\circ .
\]
We draw the frame at the mirror root, $\pb = 79.8605^\circ$, held at $79.86^\circ$ as in the worked examples of~\citep{DerrazGN2026}, for which $\sin\pb \approx 0.98438$ and $\cos\pb \approx 0.17605$; \cref{subsec:why-sign} says why this root is the natural one to draw. Both the scale unit and the root are conventions. In units of $10^{9}$~kg\,C$^{-1}$ the drawing length would be $11.6$, and the Axiom would have no root; in Planck or Stoney units it is $1$, and the Axiom returns the right angle of the complex plane. \Cref{subsec:angle-units} shows that no measured quantity depends on the choice.

\subsection{The real reading}\label{subsec:real-reading}

A product of two cosmic masses is a General Number, while an instrument returns one real number. We fix the rule once.

\begin{postulate}[The real reading]\label{def:real-reading}
Every measured force, energy and momentum is the real part of a General Number. The Composer coordinate of a force, an energy or a momentum is not read as a measured quantity.
\end{postulate}

The shadow~\cref{eq:shadow} is the other natural rule. For a neutral body the two agree; \cref{subsec:why-sign} shows where they part and why the framework reads by the real part. \Cref{tab:dictionary} sets the words of the General Numbers beside the names they take for the cosmic mass.

\begin{table}[t]
\caption{One object, two vocabularies. On the left, the words of the General Numbers, with the statement of~\citep{DerrazGN2026} that defines each one; on the right, their names for the cosmic mass in this paper.}
\label{tab:dictionary}
\small
\begin{tabularx}{\textwidth}{@{}>{\raggedright\arraybackslash}p{0.44\textwidth}c>{\raggedright\arraybackslash}X@{}}
\toprule
General Numbers & \citep{DerrazGN2026} & cosmic mass \\
\midrule
General Number $w = a + b\,\Gp$ & Ax.~3.1 & cosmic mass $\mu = m + q\,\Gp$ \\
real axis; Composer axis & Ax.~3.1 & mass axis; charge axis \\
real part $\Re(w) = a$ & Def.~3.9 & mass $m$ \\
Composer coordinate $b$ & Def.~3.9 & charge $q$ \\
Composer part $b\,\Gp$ & Def.~3.9 & charge part $q\,\Gp$, in kilograms \\
length of the Composer Unit & Ax.~3.1 & size $s = \sqrt{K/G}$; drawing length $A = s/u$ in a chosen unit $u$ \\
real number; pure Composer number & Def.~3.9 & neutral body; a charge without mass \\
real part of a product & Prop.~4.2 & every measured force, energy and momentum \\
Rule of the Square, $\Re(\Gp^2) = -\lvert\Gp\rvert^2$ & Prop.~4.1 & the minus sign between gravity and Coulomb's force, at every angle \\
mirror root & Thm~3.6 & the angle $79.86^\circ$ at which the frame is drawn, for $u = 10^{10}$~kg\,C$^{-1}$ \\
\bottomrule
\end{tabularx}
\end{table}
